\section{The title}\label{s:title}

\lem{1}{On the Meaning of Life}

\paragraph{The form of the title.} The title has the form of the classical treatise title, \emph{De} with the ablative: \emph{De brevitate vitae}, \emph{De senectute}, \emph{De finibus bonorum et malorum}. Such a title announces a subject rather than a thesis, and the text that follows it is expected to say what is true about the subject. The catalogue question to which EXI-01-a is the principal answer is phrased differently: \enquote{What is the meaning of life?} The title takes the subject of that question and drops its interrogative form; the text supplies the answer in its definition, and the remark at its end says that the answer has been given (\Lm{9}).

\paragraph{\enquote{the Meaning}.} The definite article presupposes that there is one meaning of life to be stated. On the first of the three readings distinguished in Section~\ref{s:intro}, the presupposition is substantive and contested: it presupposes that human life as such has a point, and the text would owe an account of what that point is. On the second reading, the presupposition is harmless: \enquote{the meaning of life} denotes the property in virtue of which lives are meaningful, and the text says what that property is. On the third, \enquote{the meaning} denotes the content of the concept, and the text would owe an analysis of what we mean.

The text takes the second reading, and its definition shows this. Definition~1 does not say what human life is for. It says what a life's being meaningful to a degree consists in. A theory that answered the first question would have the form \enquote{human life exists in order to $X$}, and the meaning of an individual life would derive from its contribution to $X$. Definition~1 has no such form: the projects of objective worth with which a life may be engaged are many, they need not be related to one another, and nothing in the definition requires that they serve a single end. The text therefore uses the definite article in the way in which one speaks of \emph{the} nature of knowledge or \emph{the} source of normativity: as naming whatever it is that the theory identifies, without presupposing that it is a purpose.

The second reading has a cost that the text accepts. Many people who ask what the meaning of life is are asking the first question, and an answer to the second may seem to change the subject. \citet{Metz2013} records the distinction and treats the second question as the one that analytic philosophy can address; \citet{Wolf2010} addresses the question of meaning in individual lives and does not make her answer depend on whether human life as such has a point. EXI-01-a follows both. It does not deny that the first question has a sense. It holds, and EXI-11 argues, that the answer to the second does not depend on the answer to the first: a life can be meaningful though human life as such has no cosmic purpose (Theorem~3.3 of EXI-11-a). The first question is not dismissed but separated, and EXI-06 and EXI-11 take it up: the universe does not care about us (EXI-06-a), and a cosmic purpose, if there is one, would add to the meaning available to a life without being required for it (Proposition~4.1 of EXI-11-a).

\paragraph{\enquote{of Life}.} \enquote{Life} in the title is used without an article, as a mass term, and Definition~1 replaces it with a count noun: \enquote{A life $L$}. The shift is significant. It locates the bearer of meaning in individual lives rather than in life as such, or in humanity, or in the biological phenomenon. On the first reading of the title, the bearer would be human life as a whole, and individual lives would be meaningful derivatively. On the reading the text adopts, individual lives are the primary bearers, and \enquote{the meaning of life} names a property that individual lives have to various degrees. The comment on \Lm{2.1} examines what a life is for the purposes of the definition.

\paragraph{The subject of the discipline.} The heading under which EXI-01-a stands, \enquote{Meaning, Death and the Absurd}, groups the question of meaning with two others. The grouping reflects a historical connection: in Tolstoy and in Camus, the question of meaning arises from the thought of death and is answered, or left unanswered, under the heading of the absurd. EXI-01-a severs the first connection and reshapes the second. It does not mention death, and Definition~1 makes the meaning of a life depend on what the life consists in rather than on whether it ends; EXI-19 adds that death gives a life its shape, which bears on one kind of meaning, the meaning of a life as a whole, without being a source of meaning (Proposition~1.1 of EXI-19-b). It mentions the absurd only in Remark~5, and there it treats the absurd as a challenge to one of its premises rather than as a feature of the condition its definition describes (\Lm{11}).

\paragraph{Parallels.} The title is shared, nearly or exactly, with several works in the literature: \citet{Cottingham2003}, \emph{On the Meaning of Life}, a defence of the view that meaning requires a religious framework of hope; Taylor's chapter \enquote{The Meaning of Life} \citep[ch.~18]{Taylor1970}; and the anthologies that collect the debate under that heading. EXI-01-a's answer is opposed to Cottingham's, since it makes meaning available without any religious framework, and it agrees with EXI-11 that such a framework would add to meaning without being required for it. Its answer is opposed to Taylor's, since it requires that the objects of engagement have worth that the engagement does not confer.

\paragraph{Rival readings.} Two readings of the title should be recorded and set aside. On a deflationary reading, the title is ironic: it promises an answer to the great question and delivers a definition, and the brevity of the text signals that the great question has no answer worth more than a page. Nothing in the text supports this reading. Remark~5 says plainly that the theorem answers the question of what the meaning of life is, and the register of the text is that of the rest of the discipline. On an ambitious reading, the title claims that the text answers all three questions. Remark~5 also refutes this reading, since it states two questions that the theorem does not answer (\Lm{10}, \Lm{11}). The correct reading is the one the text itself gives: the title names the subject; Definition~1 says what the meaning of life, in the sense of the second question, consists in; Theorem~4 says that it is available to every agent; and Remark~5 marks the boundaries of both claims.
